% Options for packages loaded elsewhere
\PassOptionsToPackage{unicode}{hyperref}
\PassOptionsToPackage{hyphens}{url}
%
\documentclass[
]{article}
\usepackage{amsmath,amssymb}
\usepackage{iftex}
\ifPDFTeX
  \usepackage[T1]{fontenc}
  \usepackage[utf8]{inputenc}
  \usepackage{textcomp} % provide euro and other symbols
\else % if luatex or xetex
  \usepackage{unicode-math} % this also loads fontspec
  \defaultfontfeatures{Scale=MatchLowercase}
  \defaultfontfeatures[\rmfamily]{Ligatures=TeX,Scale=1}
\fi
\usepackage{lmodern}
\ifPDFTeX\else
  % xetex/luatex font selection
\fi
% Use upquote if available, for straight quotes in verbatim environments
\IfFileExists{upquote.sty}{\usepackage{upquote}}{}
\IfFileExists{microtype.sty}{% use microtype if available
  \usepackage[]{microtype}
  \UseMicrotypeSet[protrusion]{basicmath} % disable protrusion for tt fonts
}{}
\makeatletter
\@ifundefined{KOMAClassName}{% if non-KOMA class
  \IfFileExists{parskip.sty}{%
    \usepackage{parskip}
  }{% else
    \setlength{\parindent}{0pt}
    \setlength{\parskip}{6pt plus 2pt minus 1pt}}
}{% if KOMA class
  \KOMAoptions{parskip=half}}
\makeatother
\usepackage{xcolor}
\usepackage[margin=2.5cm]{geometry}
\usepackage{color}
\usepackage{fancyvrb}
\newcommand{\VerbBar}{|}
\newcommand{\VERB}{\Verb[commandchars=\\\{\}]}
\DefineVerbatimEnvironment{Highlighting}{Verbatim}{commandchars=\\\{\}}
% Add ',fontsize=\small' for more characters per line
\newenvironment{Shaded}{}{}
\newcommand{\AlertTok}[1]{\textcolor[rgb]{1.00,0.00,0.00}{\textbf{#1}}}
\newcommand{\AnnotationTok}[1]{\textcolor[rgb]{0.38,0.63,0.69}{\textbf{\textit{#1}}}}
\newcommand{\AttributeTok}[1]{\textcolor[rgb]{0.49,0.56,0.16}{#1}}
\newcommand{\BaseNTok}[1]{\textcolor[rgb]{0.25,0.63,0.44}{#1}}
\newcommand{\BuiltInTok}[1]{\textcolor[rgb]{0.00,0.50,0.00}{#1}}
\newcommand{\CharTok}[1]{\textcolor[rgb]{0.25,0.44,0.63}{#1}}
\newcommand{\CommentTok}[1]{\textcolor[rgb]{0.38,0.63,0.69}{\textit{#1}}}
\newcommand{\CommentVarTok}[1]{\textcolor[rgb]{0.38,0.63,0.69}{\textbf{\textit{#1}}}}
\newcommand{\ConstantTok}[1]{\textcolor[rgb]{0.53,0.00,0.00}{#1}}
\newcommand{\ControlFlowTok}[1]{\textcolor[rgb]{0.00,0.44,0.13}{\textbf{#1}}}
\newcommand{\DataTypeTok}[1]{\textcolor[rgb]{0.56,0.13,0.00}{#1}}
\newcommand{\DecValTok}[1]{\textcolor[rgb]{0.25,0.63,0.44}{#1}}
\newcommand{\DocumentationTok}[1]{\textcolor[rgb]{0.73,0.13,0.13}{\textit{#1}}}
\newcommand{\ErrorTok}[1]{\textcolor[rgb]{1.00,0.00,0.00}{\textbf{#1}}}
\newcommand{\ExtensionTok}[1]{#1}
\newcommand{\FloatTok}[1]{\textcolor[rgb]{0.25,0.63,0.44}{#1}}
\newcommand{\FunctionTok}[1]{\textcolor[rgb]{0.02,0.16,0.49}{#1}}
\newcommand{\ImportTok}[1]{\textcolor[rgb]{0.00,0.50,0.00}{\textbf{#1}}}
\newcommand{\InformationTok}[1]{\textcolor[rgb]{0.38,0.63,0.69}{\textbf{\textit{#1}}}}
\newcommand{\KeywordTok}[1]{\textcolor[rgb]{0.00,0.44,0.13}{\textbf{#1}}}
\newcommand{\NormalTok}[1]{#1}
\newcommand{\OperatorTok}[1]{\textcolor[rgb]{0.40,0.40,0.40}{#1}}
\newcommand{\OtherTok}[1]{\textcolor[rgb]{0.00,0.44,0.13}{#1}}
\newcommand{\PreprocessorTok}[1]{\textcolor[rgb]{0.74,0.48,0.00}{#1}}
\newcommand{\RegionMarkerTok}[1]{#1}
\newcommand{\SpecialCharTok}[1]{\textcolor[rgb]{0.25,0.44,0.63}{#1}}
\newcommand{\SpecialStringTok}[1]{\textcolor[rgb]{0.73,0.40,0.53}{#1}}
\newcommand{\StringTok}[1]{\textcolor[rgb]{0.25,0.44,0.63}{#1}}
\newcommand{\VariableTok}[1]{\textcolor[rgb]{0.10,0.09,0.49}{#1}}
\newcommand{\VerbatimStringTok}[1]{\textcolor[rgb]{0.25,0.44,0.63}{#1}}
\newcommand{\WarningTok}[1]{\textcolor[rgb]{0.38,0.63,0.69}{\textbf{\textit{#1}}}}
\usepackage{longtable,booktabs,array}
\usepackage{calc} % for calculating minipage widths
% Correct order of tables after \paragraph or \subparagraph
\usepackage{etoolbox}
\makeatletter
\patchcmd\longtable{\par}{\if@noskipsec\mbox{}\fi\par}{}{}
\makeatother
% Allow footnotes in longtable head/foot
\IfFileExists{footnotehyper.sty}{\usepackage{footnotehyper}}{\usepackage{footnote}}
\makesavenoteenv{longtable}
\setlength{\emergencystretch}{3em} % prevent overfull lines
\providecommand{\tightlist}{%
  \setlength{\itemsep}{0pt}\setlength{\parskip}{0pt}}
\setcounter{secnumdepth}{-\maxdimen} % remove section numbering
\ifLuaTeX
\usepackage[bidi=basic]{babel}
\else
\usepackage[bidi=default]{babel}
\fi
\babelprovide[main,import]{french}
% get rid of language-specific shorthands (see #6817):
\let\LanguageShortHands\languageshorthands
\def\languageshorthands#1{}
\ifLuaTeX
  \usepackage{selnolig}  % disable illegal ligatures
\fi
\usepackage{bookmark}
\IfFileExists{xurl.sty}{\usepackage{xurl}}{} % add URL line breaks if available
\urlstyle{same}
\hypersetup{
  pdflang={fr},
  hidelinks,
  pdfcreator={LaTeX via pandoc}}

\author{}
\date{}

\begin{document}

\section{Prompt : Conception d'un TP universitaire -- Acquisition,
stockage et exploitation de données avec Google Colab et
SQL}\label{prompt-conception-dun-tp-universitaire-acquisition-stockage-et-exploitation-de-donnuxe9es-avec-google-colab-et-sql}

\subsection{1. Contexte pédagogique}\label{contexte-puxe9dagogique}

Agis en tant qu'enseignant-chercheur spécialisé en informatique, bases
de données et analyse de données.

Conçois un \textbf{travail pratique (TP) complet, progressif et
réalisable sur Google Colab}, destiné à des étudiants universitaires
découvrant l'intégration de Python, Google Drive et SQL dans un projet
d'analyse de données.

Le TP doit être organisé autour d'un \textbf{mini-projet concret}, fondé
sur l'exploitation d'un jeu de données réel.

\textbf{Principe pédagogique :} les étudiants doivent réaliser eux-mêmes
les différentes étapes. Le document doit donc présenter des consignes,
des questions, des indications et des espaces de réponse, sans fournir
immédiatement toutes les solutions.

\subsection{2. Objectifs pédagogiques}\label{objectifs-puxe9dagogiques}

À l'issue du TP, les étudiants devront être capables de :

\begin{enumerate}
\def\labelenumi{\arabic{enumi}.}
\tightlist
\item
  Rechercher et sélectionner un jeu de données public pertinent.
\item
  Télécharger un dataset depuis une source fiable.
\item
  Importer et stocker les données dans Google Drive.
\item
  Connecter Google Colab à Google Drive.
\item
  Charger et explorer un dataset avec Python et Pandas.
\item
  Créer une base de données SQLite à partir du dataset.
\item
  Exécuter des requêtes SQL depuis Python.
\item
  Réaliser des analyses statistiques simples avec SQL.
\item
  Interpréter les résultats obtenus.
\item
  Sauvegarder le notebook et les résultats du projet.
\end{enumerate}

\subsection{3. Choix du dataset}\label{choix-du-dataset}

L'étudiant doit sélectionner un jeu de données provenant de l'une des
sources suivantes :

\begin{itemize}
\tightlist
\item
  \textbf{Kaggle} : \url{https://www.kaggle.com/datasets}
\item
  \textbf{Fonds monétaire international (FMI)} :
  \url{https://www.imf.org/en/Data}
\item
  \textbf{Haut-Commissariat au Plan du Maroc (HCP)} :
  \url{https://www.hcp.ma}
\end{itemize}

Le dataset doit respecter les critères suivants :

\begin{itemize}
\tightlist
\item
  Être public et accessible légalement.
\item
  Être disponible en CSV ou Excel, ou dans un format facilement
  convertible.
\item
  Contenir au moins 100 observations et 5 variables, sauf justification.
\item
  Comporter des variables numériques et catégorielles.
\item
  Permettre l'utilisation de requêtes SQL d'agrégation.
\item
  Être suffisamment simple pour être exploité pendant un TP.
\end{itemize}

Proposer trois exemples de datasets pertinents, idéalement liés à
l'économie, à l'emploi, à la démographie ou aux entreprises.

Pour chaque dataset, indiquer la source, le lien d'accès, les
principales variables et les analyses SQL possibles. Ne pas inventer de
liens ni de caractéristiques non vérifiées.

\subsection{4. Déroulement du TP}\label{duxe9roulement-du-tp}

\subsubsection{Partie 1 -- Recherche et téléchargement du
dataset}\label{partie-1-recherche-et-tuxe9luxe9chargement-du-dataset}

Demander aux étudiants de :

\begin{enumerate}
\def\labelenumi{\arabic{enumi}.}
\tightlist
\item
  Choisir une source de données.
\item
  Identifier un dataset pertinent.
\item
  Télécharger le fichier.
\item
  Présenter brièvement son contenu.
\item
  Justifier leur choix.
\end{enumerate}

\textbf{Questions :}

\begin{itemize}
\tightlist
\item
  Quelle est la source du dataset ?
\item
  Quel est le domaine étudié ?
\item
  Combien d'observations et de variables contient-il ?
\item
  Quelles analyses pourrait-on réaliser à partir de ces données ?
\end{itemize}

\subsubsection{Partie 2 -- Stockage sur Google
Drive}\label{partie-2-stockage-sur-google-drive}

Les étudiants doivent :

\begin{enumerate}
\def\labelenumi{\arabic{enumi}.}
\tightlist
\item
  Créer un dossier \texttt{TP\_SQL\_Dataset} sur Google Drive.
\item
  Y déposer le dataset téléchargé.
\item
  Créer un notebook Google Colab.
\item
  Monter Google Drive dans Colab.
\item
  Vérifier l'existence du fichier.
\item
  Charger le dataset avec Pandas.
\end{enumerate}

Fournir uniquement le code de départ nécessaire à la connexion Google
Drive et au chargement du fichier.

\subsubsection{Partie 3 -- Exploration et préparation des
données}\label{partie-3-exploration-et-pruxe9paration-des-donnuxe9es}

Demander aux étudiants d'utiliser Pandas pour :

\begin{itemize}
\tightlist
\item
  Afficher les 10 premières lignes.
\item
  Identifier les noms et les types des colonnes.
\item
  Déterminer le nombre d'observations.
\item
  Identifier les valeurs manquantes.
\item
  Rechercher les doublons.
\item
  Calculer quelques statistiques descriptives.
\item
  Effectuer les transformations nécessaires avant l'importation dans
  SQLite.
\end{itemize}

\subsubsection{Partie 4 -- Création de la base de données
SQLite}\label{partie-4-cruxe9ation-de-la-base-de-donnuxe9es-sqlite}

Les étudiants doivent :

\begin{enumerate}
\def\labelenumi{\arabic{enumi}.}
\tightlist
\item
  Créer une base SQLite dans leur dossier Google Drive.
\item
  Créer une table adaptée à la structure du dataset.
\item
  Importer les données dans cette table.
\item
  Vérifier le nombre de lignes importées.
\item
  Afficher les premières lignes avec une requête SQL.
\end{enumerate}

Utiliser les bibliothèques :

\begin{Shaded}
\begin{Highlighting}[]
\ImportTok{import}\NormalTok{ sqlite3}
\ImportTok{import}\NormalTok{ pandas }\ImportTok{as}\NormalTok{ pd}
\end{Highlighting}
\end{Shaded}

Expliquer brièvement le rôle de \texttt{sqlite3}, de \texttt{pandas} et
de la méthode \texttt{to\_sql()}.

\subsubsection{Partie 5 -- Exploitation du dataset avec
SQL}\label{partie-5-exploitation-du-dataset-avec-sql}

Construire une série de 15 exercices SQL de difficulté progressive.

\textbf{Niveau 1 -- Requêtes fondamentales}

\begin{enumerate}
\def\labelenumi{\arabic{enumi}.}
\tightlist
\item
  Afficher toutes les observations.
\item
  Sélectionner certaines colonnes.
\item
  Limiter le nombre de résultats.
\item
  Trier les observations.
\item
  Filtrer les données avec \texttt{WHERE}.
\end{enumerate}

\textbf{Niveau 2 -- Agrégations}

\begin{enumerate}
\def\labelenumi{\arabic{enumi}.}
\setcounter{enumi}{5}
\tightlist
\item
  Calculer le nombre total d'enregistrements.
\item
  Calculer une moyenne.
\item
  Déterminer les valeurs minimale et maximale.
\item
  Regrouper les observations avec \texttt{GROUP\ BY}.
\item
  Filtrer les groupes avec \texttt{HAVING}.
\end{enumerate}

\textbf{Niveau 3 -- Analyses avancées}

\begin{enumerate}
\def\labelenumi{\arabic{enumi}.}
\setcounter{enumi}{10}
\tightlist
\item
  Combiner plusieurs conditions avec \texttt{AND} et \texttt{OR}.
\item
  Utiliser \texttt{BETWEEN} ou \texttt{IN}.
\item
  Réaliser un classement avec \texttt{ORDER\ BY}.
\item
  Construire une sous-requête SQL.
\item
  Produire une analyse synthétique répondant à une question métier ou
  économique.
\end{enumerate}

Les requêtes doivent être adaptées aux colonnes réellement présentes
dans le dataset choisi.

\subsubsection{Partie 6 -- Analyse et
interprétation}\label{partie-6-analyse-et-interpruxe9tation}

Demander aux étudiants de sélectionner cinq requêtes SQL significatives
et de :

\begin{itemize}
\tightlist
\item
  Présenter les résultats sous forme de tableaux Pandas.
\item
  Interpréter les résultats.
\item
  Produire au moins deux graphiques avec Matplotlib.
\item
  Formuler trois conclusions sur les données analysées.
\end{itemize}

\subsubsection{Partie 7 -- Livrables}\label{partie-7-livrables}

Chaque étudiant ou groupe doit remettre :

\begin{enumerate}
\def\labelenumi{\arabic{enumi}.}
\tightlist
\item
  Un notebook Google Colab fonctionnel (\texttt{.ipynb}).
\item
  Le dataset original.
\item
  La base SQLite (\texttt{.db}).
\item
  Les 15 requêtes SQL et leurs résultats.
\item
  Une synthèse d'une page présentant les principaux résultats et
  difficultés rencontrées.
\end{enumerate}

\subsection{5. Modalités
d'évaluation}\label{modalituxe9s-duxe9valuation}

Proposer une grille d'évaluation sur 20 points :

\begin{longtable}[]{@{}lr@{}}
\toprule\noalign{}
Critère & Points \\
\midrule\noalign{}
\endhead
\bottomrule\noalign{}
\endlastfoot
Choix et justification du dataset & 2 \\
Téléchargement et stockage sur Drive & 3 \\
Exploration et préparation des données & 3 \\
Création de la base SQLite & 3 \\
Qualité des requêtes SQL & 5 \\
Analyse et interprétation & 3 \\
Organisation et documentation du notebook & 1 \\
\textbf{Total} & \textbf{20} \\
\end{longtable}

\subsection{6. Contraintes techniques et
pédagogiques}\label{contraintes-techniques-et-puxe9dagogiques}

\begin{itemize}
\tightlist
\item
  Le TP doit être réalisable dans Google Colab sans installation
  complexe.
\item
  Privilégier les bibliothèques Python déjà disponibles dans Colab.
\item
  Utiliser SQLite, sans serveur de base de données externe.
\item
  Prévoir une méthode de téléchargement manuelle pour les sources
  nécessitant une authentification, notamment Kaggle.
\item
  Éviter de demander aux étudiants de publier des identifiants ou des
  clés API.
\item
  Prévoir des messages de vérification après chaque étape.
\item
  Adapter la difficulté à des étudiants débutants ou intermédiaires.
\item
  Prévoir une durée totale de 3 à 4 heures.
\item
  Distinguer clairement les consignes obligatoires et les exercices
  facultatifs.
\end{itemize}

\subsection{7. Format de sortie attendu}\label{format-de-sortie-attendu}

Générer deux livrables distincts :

\textbf{Document étudiant}

\begin{itemize}
\tightlist
\item
  Présentation du TP et objectifs.
\item
  Consignes détaillées et progressives.
\item
  Questions et exercices numérotés.
\item
  Squelettes de code à compléter.
\item
  Espaces réservés aux observations et aux réponses.
\item
  Grille d'évaluation.
\end{itemize}

\textbf{Corrigé enseignant}

\begin{itemize}
\tightlist
\item
  Notebook Google Colab complet et exécutable.
\item
  Code Python commenté.
\item
  Requêtes SQL corrigées.
\item
  Résultats attendus et explications.
\item
  Solutions alternatives lorsque cela est pertinent.
\item
  Erreurs fréquentes et pistes de remédiation.
\end{itemize}

Le corrigé doit s'appuyer sur un dataset réel, accessible et précisément
identifié. Les résultats chiffrés doivent être obtenus par exécution du
code, et non inventés.

\textbf{Résultat attendu :} un TP universitaire autonome, directement
exploitable en séance, permettant aux étudiants de suivre toute la
chaîne de traitement :

\textbf{Source de données → Téléchargement → Google Drive → Google Colab
→ Pandas → SQLite → Requêtes SQL → Analyse des résultats.}

\end{document}
